% Einheit 3 – Krümmung und Wendepunkte
\setcounter{aufgabe}{23}
\hypertarget{e3}{}\einheitenkopf{Einheit 3 von 5 · Krümmung und Wendepunkte}
\zweigzeile{Hier lernst du, mit der zweiten Ableitung die Krümmung eines Graphen und seine Wendepunkte zu bestimmen · kennst du seit Q1 (Klasse 11) · Abitur GK · baut auf: Monotonie (Einheit 1), Extrempunkte (Einheit 2), Ableitungen bilden, Gleichungen lösen, Funktionswerte berechnen}

\begin{aufgabe}{Ich kann die zweite und die dritte Ableitung bilden. Bilde $f''$ und $f'''$.}
\beispiel{f(x) &= x^4 - 2x^3 + x \\ f'(x) &= 4x^3 - 6x^2 + 1 \\ f''(x) &= 12x^2 - 12x \\ f'''(x) &= 24x - 12}
\begin{geruest}
\gz{a}{$f(x) = x^3 + 2x^2$}{\feldl{f''(x)}\feld{f'''(x)}}
\gz{b}{$f(x) = x^4 - x^2$}{\feldl{f''(x)}\feld{f'''(x)}}
\gz{c}{$f(x) = -2x^3 + 5x - 1$}{\feldl{f''(x)}\feld{f'''(x)}}
\gz{d}{$f(x) = 0{,}5x^4 + x^3$}{\feldl{f''(x)}\feld{f'''(x)}}
\end{geruest}
\end{aufgabe}

\begin{aufgabe}{Ich kann Wendepunkte berechnen. Bestimme den Wendepunkt oder die Wendepunkte des Graphen von $f$.}
\beispiel{f(x) &= x^3 + 3x^2 - 2 \\ f''(x) &= 6x + 6 = 0 && \Leftrightarrow x = -1 \\ f'''(x) &= 6 \ne 0 \\ f(-1) &= -1 + 3 - 2 = 0 && \Rightarrow W(-1 \mid 0)}
\begin{teile}
\teil $f(x) = 2x^3 - 6x^2 + 5x$
\teil $f(x) = x^4 - 6x^2 + 1$
\teil Gib für $f$ aus b) an, wo der Graph linksgekrümmt und wo er rechtsgekrümmt ist.
\teil Bestimme den Wendepunkt von $f(x) = x^3 - 6x^2 + 11x - 5$ und die Gleichung der Wendetangente.
\teil Weise nach, dass $W(1 \mid 3)$ ein Wendepunkt des Graphen von $f(x) = -x^3 + 3x^2 + 1$ ist.
\teil Zeige, dass der Graph von $f(x) = 3x^4 + 4x^3 + 1$ im Punkt $S(0 \mid 1)$ einen Sattelpunkt hat.
\end{teile}
\end{aufgabe}

\begin{aufgabe}{Ich kann Wendepunkte berechnen – weiter.}
\begin{teile}
\teil Gegeben ist $f(x) = x\cdot\mathrm{e}^{x}$. Zeige, dass $f''(x) = (x + 2)\cdot\mathrm{e}^{x}$ gilt, und bestimme den Wendepunkt.
\teil Eine Funktion $g$ ist auf $\mathbb{R}$ definiert. Ihr Graph nähert sich für $x \to -\infty$ von unten der $x$-Achse, hat genau einen Extrempunkt, den Tiefpunkt $T(1 \mid -3)$, und steigt für $x > 1$ unbeschränkt. Skizziere einen möglichen Graphen und begründe ohne Rechnung, dass der Graph von $g$ einen Wendepunkt hat. (Abitur 2021)
\end{teile}

\begin{ksys}[xmin=-6,xmax=4,ymin=-4,ymax=3]
\end{ksys}

\end{aufgabe}

\begin{aufgabe}{Ich kann Geraden durch Wendepunkte untersuchen.}
\begin{teile}
\teil Berechne für die Wendetangente aus Aufgabe 25 d) den Winkel, den sie mit der $x$-Achse einschließt.
\end{teile}
Die Abbildung zeigt den Graphen von $f(x) = x^4 - 2x^3$ für $-1 \le x \le 2{,}2$.

\begin{ksys}[xmin=-2,xmax=3,ymin=-2,ymax=4]
\funktionab{\x^4-2*\x^3}{f}{-1}{2.2}
\end{ksys}

\begin{teile}
\teil Berechne die beiden Wendepunkte, stelle die Gleichung der Geraden durch beide Wendepunkte auf und zeichne sie ein.
\teil Zeichne eine zu dieser Geraden parallele Gerade ein, die mit dem abgebildeten Graphen genau einen gemeinsamen Punkt hat, und gib ihre Gleichung an.
\end{teile}
\end{aufgabe}

\begin{aufgabe}{Ich finde den Fehler beim Krümmungsverhalten. Lea untersucht die Krümmung des Graphen von $f(x) = -x^3 + 6x^2$:}
\rechnung{f''(x) &= -6x + 12 = 0 && \Leftrightarrow x = 2 \\ f''(0) &= 12 > 0 && \Rightarrow \text{rechtsgekrümmt auf } ]{-\infty};\,2] \\ f''(3) &= -6 < 0 && \Rightarrow \text{linksgekrümmt auf } [2;\,\infty[}
\begin{teile}
\teil Finde den Fehler, benenne ihn und korrigiere die Rechnung.
\teil Bestimme, wo der Graph von $g(x) = x^3 + 3x^2 - x$ linksgekrümmt und wo er rechtsgekrümmt ist.
\end{teile}
\end{aufgabe}

\begin{aufgabe}{Ich kann begründen, was an einem Wendepunkt geschieht.}
\begin{teile}
\teil Begründe, warum $f''$ an einer Stelle das Vorzeichen wechseln muss, damit dort ein Wendepunkt liegt.
\teil Begründe, warum die Wendetangente den Graphen im Wendepunkt durchsetzt, also auf die andere Seite des Graphen wechselt.
\teil Untersuche, für welche Steigungen $m$ die Gerade $y = mx$ durch den Wendepunkt $W(0 \mid 0)$ den Graphen von $f(x) = x^3 + x$ in genau einem Punkt und für welche in drei Punkten schneidet.
\end{teile}
\end{aufgabe}

\begin{aufgabe}{Ich kann die Krümmung im Sachzusammenhang deuten. Ein Straßenstück wird von oben gesehen durch $f(x) = 0{,}1x^3 - 0{,}6x^2 + 2$ für $0 \le x \le 5$ beschrieben ($x$ und $y$ in 100 m). Ein Auto fährt in Richtung wachsender $x$-Werte.}
\begin{teile}
\teil Bestimme den Punkt, an dem das Auto von einer Rechtskurve in eine Linkskurve wechselt.
\teil Gib an, auf welchem Abschnitt der Fahrer nach links lenkt.
\end{teile}
\end{aufgabe}
